% =====================================================================
%  Section 5 -- Results
% =====================================================================

\section{Results}
\label{sec:results}

This section is organised around Tables~\ref{tab:t1-footprints}--\ref{tab:t10-combined} and Figures~\ref{fig:f2-pareto}--\ref{fig:f4-throughput}. The pipeline schematic (Figure~\ref{fig:f1-pipeline}) is shown in Section~\ref{sec:methods} and the qualitative panel (Figure~\ref{fig:f3-qualitative}) in Section~\ref{sec:results-pipeline}.

\subsection{Teacher Performance and Karc{\i} \emph{et al.}~(2026) Comparison}
\label{sec:results-teacher}

Our UNet $+$ EfficientNet-B0 teacher achieves \textbf{IoU $= \teacherIoU \pm \teacherStd$} on FSSD across the three folds (F1 $= 0.9611 \pm 0.0044$, accuracy $= 0.9729 \pm 0.0025$). The reference UNet $+$ EfficientNet-B7 number reported by Karc{\i} \emph{et al.}~\cite{karci2026deep} is IoU $= 0.927$, achieved with approximately $10\times$ the parameters ($\sim 63$\,M vs.~our 6.25\,M). The 0.0013 IoU gap between our B0 teacher and the published B7 reference is within the cross-fold standard deviation of both studies and indicates that the FSSD benchmark saturates well before the B7 capacity is required; B0 is sufficient as the teacher.

\subsection{Students Match or Exceed the Teacher Without Distillation (Negative-KD Finding)}
\label{sec:results-kd}

The full segmentation results across all 14 configurations $\times$ 3 folds appear in Table~\ref{tab:t2-segmentation}. Three observations:

\input{tables/T2_segmentation}

\begin{enumerate}
\item \textbf{The best student exceeds the teacher.} MobileViT-XXS in the no-KD configuration attains IoU $= \bestStudentIoU \pm \bestStudentStd$, \emph{0.011 IoU points above the teacher} (\teacherIoU $\pm$ \teacherStd). The student is also $2.0\times$ smaller in parameters (3.09\,M vs.~6.25\,M). This is the headline result against which the original KD-recipe framing of this paper has to be evaluated.
\item \textbf{EfficientNet-Lite0 (no-KD) matches the teacher within standard deviation} ($0.9277 \pm 0.0047$ vs.~$\teacherIoU \pm \teacherStd$). MobileNetV3-Small (no-KD) likewise matches it ($0.9269 \pm 0.0043$).
\item \textbf{The MobileNetV2 baseline also matches the teacher} ($0.9257 \pm 0.0032$). The teacher's parameter count is therefore not the source of any advantage; it is an over-parameterised lightweight UNet, not a ``heavy'' model in the regime where KD typically helps.
\end{enumerate}

We interpret these observations as direct evidence that FSSD, at 663 RGB images and IoU $\approx 0.93$, is saturated for the ImageNet-pretrained UNet family. The variance across $(\text{model} \times \text{KD})$ cells of Table~\ref{tab:t2-segmentation} is dominated by fold-level seed noise rather than by genuine architectural separations. KD has no headroom to transfer.

Table~\ref{tab:t3-kd-ablation} isolates the KD ablation on the best student. None of the three KD configurations improves on the no-KD baseline:

% =====================================================================
%  Table T3 -- KD Ablation on Best Student
%  Source of truth: code/runs/3fold/summary.json + stats.json (Wilcoxon)
% =====================================================================
\begin{table}[ht]
\centering
\caption{\textbf{KD ablation on the best student (MobileViT-XXS).} Wilcoxon two-sided; $n=3$ folds gives a minimum reportable $p \approx 0.25$. $\Delta$IoU is relative to no KD.}
\label{tab:t3-kd-ablation}
\small
\setlength{\tabcolsep}{4pt}
\begin{tabular}{lccc}
\toprule
\textbf{KD} & \textbf{IoU (mean $\pm$ std)} & \textbf{$\Delta$IoU} & \textbf{Wilcoxon $p$} \\
\midrule
none      & 0.9367 $\pm$ 0.0055 & ---      & --- \\
resp      & 0.9326 $\pm$ 0.0060 & $-0.0041$ & n/a \\
feat      & 0.9357 $\pm$ 0.0046 & $-0.0011$ & n/a \\
comb      & 0.9318 $\pm$ 0.0070 & $-0.0049$ & 0.250 \\
\bottomrule
\end{tabular}
\end{table}


The Wilcoxon test on no-KD vs.~comb-KD yields $p = 0.250$, the minimum two-sided p-value reportable at $n=3$, which means the fold-level test is consistent with both ``comb-KD is no worse'' and ``comb-KD is meaningfully worse'' under low power, but in no case does it support ``comb-KD is better''. Across \emph{all three architectures $\times$ three KD configurations $\times$ three folds}, the combined KD configuration produces a higher mean IoU than the no-KD baseline in zero out of three architectures.

Beyond the fold-level means, the \emph{direction} of the effect is consistent: for the best student, MobileViT-XXS, the combined-KD IoU is below the no-KD IoU in every one of the three folds, and the no-KD MobileViT-XXS exceeds the teacher in every fold (fold-wise margins $+0.013$, $+0.016$, $+0.004$ IoU; mean $+0.011$).

The per-image paired bootstrap (Section~\ref{sec:setup-stats}; $n=399$ held-out validation images, $10{,}000$ replicates; archived as \texttt{runs/3fold/stats\_bootstrap.json}) sharpens this into a powered statement. Across all nine (architecture $\times$ KD-configuration) contrasts, \emph{no} KD configuration significantly improves over its no-KD baseline: no 95\% confidence interval on $\Delta\mathrm{IoU}_{\text{no-KD}-\text{KD}}$ lies entirely below zero. Two contrasts show KD significantly \emph{worse} after Holm correction for the nine comparisons (one-sided bootstrap $p$): response-only KD on MobileViT-XXS ($\Delta = +0.0065$, 95\% CI $[+0.0029,+0.0103]$, $p=0.0001$) and combined KD on MobileViT-XXS ($+0.0076$, $[+0.0035,+0.0118]$, $p=0.0002$). Response-only KD on MobileNetV3-Small ($+0.0047$, $[+0.0008,+0.0091]$, $p=0.010$) and on EfficientNet-Lite0 ($+0.0040$, $[-0.0007,+0.0086]$, $p=0.046$) point the same way but do not survive the correction; the remaining five contrasts show no effect. The response-based term is therefore harmful for MobileViT-XXS and borderline for the other two, consistent with its interpretation as a noise-injection regulariser on a saturated benchmark (Section~\ref{sec:discussion-negative-kd}). The headline result is likewise significant: MobileViT-XXS (no-KD) exceeds the teacher by $\Delta\mathrm{IoU} = +0.015$ per image (95\% CI $[+0.010,+0.020]$, $p < 0.001$).

\subsection{INT8 Quantization is Architecture-Dependent}
\label{sec:results-quant}

Table~\ref{tab:t4-quant} summarises post-training INT8 quantization for the three deployable students:

% =====================================================================
%  Table T4 -- Quantization
%  Source of truth: code/runs/3fold/quantized.json + ONNX file sizes.
% =====================================================================
\begin{table*}[ht]
\centering
\caption{\textbf{Post-training INT8 quantization} (static QDQ, Percentile-99.999\% calibration, 80 train images). EfficientNet-Lite0 (highlighted) is the recommended deployable configuration; the other two students suffer architecture-driven PTQ collapse (Section~\ref{sec:results-quant}). For the two collapsing students the INT8 IoU is strongly multi-modal across folds (MobileNetV3-Small: $0.772 / 0.0001 / 0.657$; MobileViT-XXS: $0.718 / 0.464 / 0.244$), so their mean\,$\pm$\,std is an instability indicator, not a central tendency.}
\label{tab:t4-quant}
\small
\setlength{\tabcolsep}{4pt}
\begin{tabular}{lcccccc}
\toprule
& & & & \multicolumn{2}{c}{\textbf{ONNX size (MB)}} & \\
\cmidrule(lr){5-6}
\textbf{Model} & \textbf{FP32 IoU} & \textbf{INT8 IoU} & \textbf{$\Delta$ IoU} & \textbf{FP32} & \textbf{INT8} & \textbf{Reduction} \\
\midrule
MobileNetV3-Small   & 0.9269 $\pm$ 0.0043 & 0.4762 $\pm$ 0.3399 & $-0.4507$ & 13.8 & 3.7 & $3.69\times$ \\
\textbf{EfficientNet-Lite0} & \textbf{0.9277 $\pm$ 0.0047} & \textbf{0.8851 $\pm$ 0.0149} & $\boldsymbol{-0.0426}$ & \textbf{19.8} & \textbf{5.2} & $\boldsymbol{3.83\times}$ \\
MobileViT-XXS       & 0.9367 $\pm$ 0.0055 & 0.4752 $\pm$ 0.1936 & $-0.4616$ & 12.0 & 3.5 & $3.42\times$ \\
\bottomrule
\end{tabular}
\end{table*}


The cross-architecture spread is the central finding of this section. Under an identical Percentile-99.999\% static QDQ recipe, EfficientNet-Lite0 retains 95\% of its FP32 IoU with a tight $\pm 0.015$ cross-fold standard deviation, whereas MobileNetV3-Small and MobileViT-XXS both lose $\sim 0.45$ IoU points with cross-fold standard deviations of 0.19--0.34, an order of magnitude larger than EfficientNet-Lite0's. We interpret these failures architecturally:

\begin{itemize}
\item \textbf{MobileNetV3-Small} uses the HardSwish activation $x \cdot \mathrm{ReLU6}(x+3)/6$, which has an unbounded positive tail and produces calibration statistics dominated by single-image outliers. We verify this mechanism directly (\texttt{scripts/ptq\_activation\_diagnostic.py}; averaged over the three folds' calibration sets): across the encoder's activation layers, MobileNetV3-Small's post-activation magnitudes reach a maximum of $156.7 \pm 8.9$ (99.999-percentile $83.7$; median per-layer maximum $15.6$), whereas EfficientNet-Lite0's ReLU6 encoder is bounded at a median per-layer maximum of \emph{exactly} $6.00$ (99.999-percentile $7.0$, maximum $9.3$), a $12$--$17\times$ tighter dynamic range. Per-tensor INT8 calibration cannot fit a single scale that simultaneously preserves the activations of common flood-pixel patches and of the bright reflective outliers (wet asphalt, water-surface sun-glare) that drive HardSwish's tail. Percentile-99.999\% mitigates but does not eliminate this; we conjecture that MobileNetV3-S requires either (a)~QAT to learn quantizer-friendly activations or (b)~a HardSwish\,$\rightarrow$\,ReLU6 surgical replacement before quantization.
\item \textbf{MobileViT-XXS} combines a CNN trunk with self-attention blocks containing LayerNorm. The LayerNorm $\gamma / \beta$ tensors are rank-1 across channels, which static per-tensor calibration handles poorly; restricting quantization to \texttt{Conv/Gemm/MatMul} (Section~\ref{sec:methods-ptq}) keeps LayerNorm in FP32 but leaves the attention \texttt{MatMul} itself quantized with a single per-tensor scale, where the high dynamic range of attention scores exceeds what static calibration can fit.
\item \textbf{EfficientNet-Lite0} uses ReLU6 throughout, which keeps the encoder's activation range narrow (median per-layer maximum 6.0, overall maximum 9.3), so per-tensor activation calibration finds a stable scale in every fold. Its 0.043 IoU drop is larger than the roughly one-point accuracy loss typical of INT8 classification models, but it is stable across folds.
\end{itemize}

The practical implication: \emph{bounded-activation backbones (ReLU6) are the safe default for per-tensor static INT8 PTQ on flood segmentation; HardSwish and attention-based backbones fail under this recipe and should stay FP32 unless a stronger recipe is validated.} Our recipe uses per-tensor weight scales throughout, which are themselves a known cause of PTQ loss in MobileNet-family models~\cite{krishnamoorthi2018quantizing}; we have not tested per-channel weights, activation surgery or QAT, and any of them may recover the two failing backbones. We do not attempt QAT in this paper; it is out of scope and an order-of-magnitude larger training investment. We instead recommend EfficientNet-Lite0 as the deployable configuration on the basis of its joint accuracy $+$ size $+$ stability behaviour.

\subsection{Per-Platform Inference Latency}
\label{sec:results-latency}

Table~\ref{tab:t5-latency} reports inference latency on the two deployment platforms.

% =====================================================================
%  Table T5 -- Per-Platform Latency
%  Source of truth: code/runs/3fold/{lat_m2,lat_wasm}.json
% =====================================================================
\begin{table*}[ht]
\centering
\caption{\textbf{Per-platform inference latency.} Single-image batches, $256{\times}256$ input. M2~Pro and browser WASM both consume the same \texttt{.onnx} artifact under different execution providers. INT8 WASM rows are n/a: our INT8 QDQ artifacts require ONNX Runtime $\geq 1.18$, and ONNX Runtime Web 1.17.3 failed to load them. \textsuperscript{a}Teacher timed under the same ONNX Runtime CPU execution provider as the students (37.6\,ms p50). Under PyTorch eager CPU the same teacher is 274\,ms ($7.3\times$ slower); reporting that PyTorch figure against ORT-CPU students would inflate the apparent teacher-vs-student speedup. On a common ORT-CPU backend the FP32 students are only $1.1$--$1.9\times$ faster than the teacher, and the recommended EfficientNet-Lite0 INT8 is $2.2\times$ faster. All timing rows trace to \texttt{lat\_m2.json} (student rows via \texttt{bench\_m2.py}; teacher ORT-CPU row via \texttt{bench\_teacher\_ort.py}); see \texttt{runs/3fold/PROVENANCE.md}.}
\label{tab:t5-latency}
\small
\setlength{\tabcolsep}{5pt}
\begin{tabular}{lcccccc}
\toprule
\multirow{2}{*}{\textbf{Model}}
  & \multicolumn{3}{c}{\textbf{Apple M2~Pro (ORT-CPU)}}
  & \multicolumn{3}{c}{\textbf{Browser WASM SIMD (ORT-Web)}} \\
\cmidrule(lr){2-4} \cmidrule(lr){5-7}
  & p50 (ms) & p95 (ms) & FPS  & p50 (ms) & p95 (ms) & FPS  \\
\midrule
Teacher (UNet$+$EffB0) FP32\textsuperscript{a} & 37.6 & 44.2 & 26.6 & \multicolumn{3}{c}{n/a (out-of-scope for browser deploy)} \\
\midrule
MobileNetV3-Small FP32             &  19.5 &  20.4 & 51.2  & 133.3 & 134.8 & 7.5 \\
MobileNetV3-Small INT8             &  13.7 &  14.1 & 72.9  & \multicolumn{3}{c}{n/a (ORT-Web 1.17 limitation)} \\
EfficientNet-Lite0 FP32            &  34.2 &  52.7 & 29.3  & 163.8 & 210.5 & 6.1 \\
\textbf{EfficientNet-Lite0 INT8}   &  \textbf{17.3} & \textbf{20.1} & \textbf{57.7} & \multicolumn{3}{c}{n/a (ORT-Web 1.17 limitation)} \\
MobileViT-XXS FP32                 &  28.3 &  32.8 & 35.4  & 149.4 & 159.8 & 6.7 \\
MobileViT-XXS INT8                 &  21.6 &  25.0 & 46.3  & \multicolumn{3}{c}{n/a (ORT-Web 1.17 limitation)} \\
\bottomrule
\end{tabular}
\end{table*}


Three remarks. First, the teacher-versus-student latency comparison must be read on a common runtime. Measured under the ONNX Runtime CPU backend, the teacher runs at 37.6\,ms p50 (26.6\,FPS; ORT 1.17, re-timed at 36.7\,ms under the students' ORT 1.22.1), not the 274\,ms of PyTorch eager mode; the PyTorch\,$\rightarrow$\,ORT switch alone is a $7.3\times$ speedup with no architecture change. On this common runtime the FP32 students are only $1.1$--$1.9\times$ faster than the teacher, so the deployment advantage of the recommended EfficientNet-Lite0 comes primarily from INT8 precision (a $2.2\times$ speedup over the teacher, 17.3 vs 37.6\,ms) and from its smaller, quantization-stable binary, not from an FP32 latency gap. Second, \emph{INT8 yields a further $1.3$--$2.0\times$ speedup on top of the FP32 student baseline}, consistent with INT8 SIMD-CPU expectations. Third, \emph{the WASM column is roughly $5\times$ slower than the M2-CPU column for the same FP32 model}, which is the expected single-threaded-WebAssembly-vs-native-CPU gap and represents an effective lower bound on browser-deployment performance. INT8 WASM rows are marked n/a because ONNX Runtime Web 1.17.3 could not load the INT8 artifacts (Section~\ref{sec:methods-deploy}).

\subsection{Pareto Analysis}
\label{sec:results-pareto}

\begin{figure*}[ht]
\centering
\includegraphics[width=0.75\linewidth]{figures/F5_bubble.pdf}
\caption{\textbf{Consolidated accuracy--latency--size--precision trade-off.} Each configuration is placed by Apple~M2~Pro ORT-CPU p50 latency ($x$) and IoU ($y$); bubble area is proportional to parameter count and colour encodes ONNX disk size, so INT8 quantization moves a model left (faster) and toward a smaller, darker marker. Circles are FP32, diamonds INT8. The two red-ringed points are the recommended operating points: \textbf{EfficientNet-Lite0 INT8} (maximum throughput at acceptable accuracy) and \textbf{MobileViT-XXS FP32} (maximum accuracy); the two INT8 diamonds at IoU\,$\approx 0.48$ are MobileNetV3-Small and MobileViT-XXS, whose static-QDQ INT8 collapses (Table~\ref{tab:t4-quant}). The teacher (yellow, 22\,MB) offers no M2 latency advantage over the FP32 students. Browser WASM latencies are in Table~\ref{tab:t5-latency}; the four single-axis Pareto panels of the earlier draft remain in the repository as \texttt{figures/F2\_pareto.pdf}.}
\label{fig:f2-pareto}
\end{figure*}

This consolidated view identifies two distinct deployment operating points:

\begin{itemize}
\item \textbf{Maximum-accuracy operating point:} MobileViT-XXS no-KD FP32 (IoU 0.937, 28\,ms M2, 149\,ms WASM, 12\,MB).
\item \textbf{Maximum-throughput-with-acceptable-accuracy operating point:} EfficientNet-Lite0 INT8 (IoU 0.885, 17\,ms M2, no WASM, 5.2\,MB).
\end{itemize}

We recommend the latter for unattended flood-mapping pipelines ($1.6\times$ throughput, $2.3\times$ smaller binary) and the former for human-in-the-loop applications where the 5.2-point IoU drop matters.

\begin{figure*}[ht]
\centering
\includegraphics[width=0.92\linewidth]{figures/F4_throughput.png}
\caption{\textbf{Per-platform throughput, FP32 vs.~INT8.} Frames-per-second on M2~Pro (left) and browser WASM SIMD (right); INT8 entries are unsupported on ORT-Web 1.17.}
\label{fig:f4-throughput}
\end{figure*}

\subsection{Qualitative Examples}
\label{sec:results-pipeline}

The pipeline schematic appears at the top of Section~\ref{sec:methods} (Figure~\ref{fig:f1-pipeline}). The qualitative comparison between teacher and the deployable student is in Figure~\ref{fig:f3-qualitative}.

\begin{figure*}[ht]
\centering
\includegraphics[width=0.92\linewidth]{figures/F3_qualitative.pdf}
\caption{\textbf{Qualitative comparison on FSSD fold-0 validation.} Two rows $\times$ five columns: input, ground truth, teacher (UNet$+$EfficientNet-B0, FP32) prediction, EfficientNet-Lite0-UNet INT8 prediction, and error map (TP green, FP red, FN blue). \textbf{Easy row} (top): the highest-IoU validation image under the deployable INT8 student (FSS1604, IoU $=0.995$ for both teacher and student). \textbf{Hard row} (bottom): the lowest-IoU image with $\geq 2\%$ flood-pixel coverage (FSS904, teacher IoU $=0.546$; student IoU $=0.157$). The hard row shows the two persistent failure modes discussed in Section~\ref{sec:discussion-failures}: false-positive activations on bright reflective non-water surfaces (red regions), and false-negative under-segmentation along vegetation-occluded boundaries (blue regions).}
\label{fig:f3-qualitative}
\end{figure*}

Two persistent failure modes shared by teacher and student are visible in the hard row of Figure~\ref{fig:f3-qualitative}: wet asphalt confused with water, and under-segmentation of vegetation-occluded flood edges. We return to these in Section~\ref{sec:discussion-failures}.

\subsection{Out-of-Distribution Generalisation (Sen1Floods11)}
\label{sec:results-ood}

To test whether the FSSD-trained configurations transfer beyond their training distribution, we evaluate every model, with no fine-tuning, on 250 hand-labelled RGB chips from Sen1Floods11~\cite{bonafilia2020sen1floods11}, a Sentinel-2 optical flood benchmark whose top-down satellite viewpoint and sensor characteristics differ sharply from FSSD's close-range aerial, oblique, and ground photographs. Table~\ref{tab:t6-ood} reports the result, averaged over the three folds' checkpoints.

\input{tables/T6_ood}

The transfer is poor across the board: every model collapses from IoU $\approx 0.93$ in-distribution to IoU $\approx 0.16$ on Sen1Floods11, a near-total loss of flood-pixel overlap. Two points matter for the deployment story. First, the collapse is essentially uniform across the teacher (0.160), the MobileNetV2 baseline (0.164), and the three lightweight students (0.155--0.173): it is a domain-shift effect, not a compression or capacity artefact. The recommended EfficientNet-Lite0 ($0.155 \pm 0.013$) is within one cross-fold standard deviation of the teacher ($0.160 \pm 0.007$) on the OOD set, so the OOD collapse does not single out the deployable student relative to the teacher; the two students that edge above the teacher (MobileNetV3-Small, 0.173; MobileViT-XXS, 0.169) do so within the same cross-fold noise. Second, in absolute terms IoU $\approx 0.16$ is not operationally useful: an FSSD-trained close-range-RGB model is \emph{not} a drop-in flood mapper for Sentinel-2 satellite optical imagery, and cross-sensor deployment requires domain adaptation or retraining on in-sensor data. The precision/recall split also mirrors the operational error asymmetry (Section~\ref{sec:discussion-asymmetry}): MobileViT-XXS reaches the highest OOD recall ($0.40$) at the cost of precision (it over-predicts flood, the safer direction for event-time mapping; Section~\ref{sec:discussion-asymmetry}), whereas the MobileNetV2 baseline is the most precision-favouring ($0.385$ precision, $0.230$ recall).

\subsection{Within-Regime Zero-Shot Generalization and Quantization under Shift}
\label{sec:results-zeroshot}

The Sen1Floods11 result bounds the recommendations across \emph{sensors}; it leaves open how the models behave on unseen imagery from the regime they were built for: close-range RGB flood photography. Table~\ref{tab:t9-zeroshot} therefore evaluates the FSSD-trained models, without any fine-tuning, on two within-regime external sets with verified zero image overlap against FSSD (Section~\ref{sec:methods-data}): AIFloodSense (94 ground-level flood photographs) and the held-out FloodNet test split (424 UAV images, 215 of which contain no water). Each configuration is evaluated with all three folds' checkpoints; we report mean $\pm$ cross-fold standard deviation. Three results follow.

% =====================================================================
%  Table T9 -- Zero-shot generalization (FSSD-trained, no fine-tuning)
%  Source of truth: code/runs/3fold/zero_shot_3fold.json
%  (per-fold detail: zero_shot_fold{0,1,2}.json; aggregate_zero_shot.py)
% =====================================================================
\begin{table*}[ht]
\centering
\caption{\textbf{Zero-shot generalization of FSSD-trained models} (no fine-tuning; mean $\pm$ std across the three folds' checkpoints) on AIFloodSense (94 images, all flood-containing) and the held-out FloodNet test split (424 images, 215 of them water-free). Per-image macro metrics (Section~\ref{sec:setup-metrics}): IoU over all images, flood-only IoU over flood-containing images, and false-alarm rate (FAR) over water-free images. FSSD contains no water-free images, and every FP32 model false-alarms on 72--100\% of water-free FloodNet scenes. INT8 degradation exceeds its in-domain counterpart (cf.~Table~\ref{tab:t4-quant}): MobileNetV3-Small-INT8's seemingly moderate overall FloodNet IoU is an artifact of near-constant background prediction (flood-only IoU 0.08).}
\label{tab:t9-zeroshot}
\small
\setlength{\tabcolsep}{4pt}
\begin{tabular}{llcccc}
\toprule
& & \textbf{AIFloodSense} & \multicolumn{3}{c}{\textbf{FloodNet test}} \\
\cmidrule(lr){4-6}
\textbf{Model} & \textbf{Prec.} & \textbf{IoU} & \textbf{IoU (all)} & \textbf{IoU (flood)} & \textbf{FAR} \\
\midrule
Teacher (UNet$+$EffB0)  & FP32 & 0.706 $\pm$ 0.004 & 0.288 $\pm$ 0.014 & 0.509 $\pm$ 0.010 & 0.927 $\pm$ 0.018 \\
MobileNetV2 baseline    & FP32 & 0.690 $\pm$ 0.014 & 0.378 $\pm$ 0.015 & 0.475 $\pm$ 0.020 & 0.716 $\pm$ 0.040 \\
MobileNetV3-Small       & FP32 & 0.698 $\pm$ 0.008 & 0.285 $\pm$ 0.005 & 0.506 $\pm$ 0.012 & 0.929 $\pm$ 0.021 \\
EfficientNet-Lite0      & FP32 & 0.702 $\pm$ 0.006 & 0.388 $\pm$ 0.025 & 0.507 $\pm$ 0.021 & 0.727 $\pm$ 0.032 \\
MobileViT-XXS           & FP32 & 0.723 $\pm$ 0.006 & 0.253 $\pm$ 0.004 & 0.512 $\pm$ 0.007 & 1.000 $\pm$ 0.000 \\
\midrule
MobileNetV3-Small       & INT8 & 0.192 $\pm$ 0.164 & 0.355 $\pm$ 0.140 & 0.077 $\pm$ 0.073 & 0.375 $\pm$ 0.346 \\
EfficientNet-Lite0      & INT8 & 0.649 $\pm$ 0.028 & 0.258 $\pm$ 0.011 & 0.453 $\pm$ 0.032 & 0.932 $\pm$ 0.012 \\
MobileViT-XXS           & INT8 & 0.337 $\pm$ 0.078 & 0.092 $\pm$ 0.019 & 0.172 $\pm$ 0.049 & 0.986 $\pm$ 0.011 \\
\bottomrule
\end{tabular}
\end{table*}


\emph{First, FP32 students transfer as well as the teacher.} On AIFloodSense every FP32 model lands within 0.69--0.72 IoU with cross-fold standard deviations below 0.015; the best transfer is MobileViT-XXS ($0.723 \pm 0.006$), and the teacher holds no advantage ($0.706 \pm 0.004$). Combined KD transfers slightly \emph{worse} than no-KD on all three architectures ($-0.005$ to $-0.015$ IoU), extending the negative distillation result of Section~\ref{sec:results-kd} beyond the training distribution. On FloodNet's flood-containing images, all FP32 models cluster at 0.47--0.51 flood-IoU, a uniform $\sim$0.42-point drop from in-domain performance that reflects the viewpoint and resolution gap of high-altitude UAV imagery rather than any architecture-specific failure.

\emph{Second, FSSD-trained models cannot say ``no flood''.} Every FP32 model, teacher included, predicts water on 72--100\% of the 215 water-free FloodNet scenes (Table~\ref{tab:t9-zeroshot}, FAR column; MobileViT false-alarms on all of them in every fold). The cause is structural: FSSD contains flood imagery only, so the decision ``this scene contains no water'' is never supervised. For the monitoring use-cases that motivate edge deployment (continuous observation of scenes that are dry most of the time), this failure mode is disqualifying independently of segmentation quality, it is invisible to every FSSD-internal metric in Tables~\ref{tab:t2-segmentation}--\ref{tab:t4-quant}, and in unattended use it is the costly error (Section~\ref{sec:discussion-asymmetry}). Retraining on negative-image-bearing data (Section~\ref{sec:methods-data}) is the direct remedy; Section~\ref{sec:results-combined} tests it.

\emph{Third, INT8 degradation grows with distribution shift.} In-domain, EfficientNet-Lite0-INT8 loses 0.043 IoU (Table~\ref{tab:t4-quant}); zero-shot it loses 0.053 on AIFloodSense ($0.702 \rightarrow 0.649$) and 0.054 flood-IoU on FloodNet ($0.507 \rightarrow 0.453$), while its false-alarm rate jumps from 0.727 to 0.932: the quantized model fails the water-free scenes that its FP32 parent still partially rejects. The fragile architectures collapse outright: MobileNetV3-Small-INT8 drops to $0.192 \pm 0.164$ IoU on AIFloodSense and $0.077 \pm 0.073$ flood-IoU on FloodNet (its seemingly moderate overall FloodNet score ($0.355$) comes from predicting almost nothing, which is rewarded on the 215 water-free images), and MobileViT-XXS-INT8 loses more than half its FP32 transfer ($0.723 \rightarrow 0.337$). Per-tensor PTQ calibrated on 80 in-domain images therefore not only fails architecture-dependently (Section~\ref{sec:results-quant}) but degrades further under the distribution shift that defines real deployments, most visibly in false-alarm behaviour; we flag this as a practitioner trap that in-domain PTQ validation does not reveal, and note that it also illustrates why mixed test sets must be reported with split metrics rather than a single macro IoU.

\subsection{Retraining on Negative-Image-Bearing Data}
\label{sec:results-combined}

Section~\ref{sec:results-zeroshot} attributes the false-alarm failure to FSSD's lack of water-free images. To test that attribution, and to ask whether the negative-KD result survives on a less saturated benchmark, we retrain the teacher and the three students under no-KD and combined KD on the combined corpus of Section~\ref{sec:methods-data}. Each fold trains on 2{,}057 images: 1{,}825 FloodNet train$+$val images (855 of them water-free) and the 232 Kaggle Flood Area Segmentation images outside the held-out fold. Model selection uses the held-out Kaggle fold (54--55 images after the duplicate filter, all flood-containing), and a source-balanced sampler draws each source equally often in expectation, because FloodNet outnumbers the Kaggle images about 8:1 and carries far less water per image. Everything else follows Section~\ref{sec:setup-impl}, except that students also use batch size 8. The FloodNet test split and AIFloodSense remain untouched by training and model selection, so the two regimes are compared on identical held-out images (Table~\ref{tab:t10-combined}). To fit the free-tier budget this experiment runs the reduced grid only: no response- or feature-only KD, no MobileNetV2 baseline, and no INT8 re-evaluation.

% =====================================================================
%  Table T10 -- Combined-dataset retraining (D2) vs FSSD-trained zero-shot
%  Source of truth (combined-trained): code/runs/combined/summary.json
%    (per-fold detail: results_fold{0,1,2}.json; aggregate_combined.py)
%  Source of truth (FSSD-trained):     code/runs/3fold/zero_shot_3fold.json
%    (same rows as Table T9; std there)
% =====================================================================
\begin{table*}[ht]
\centering
\caption{\textbf{Retraining on negative-image-bearing data removes the false-alarm failure.} Left: FSSD-trained models evaluated zero-shot (means of Table~\ref{tab:t9-zeroshot}; standard deviations there). Right: the same configurations retrained on the combined corpus (Kaggle Flood Area Segmentation $+$ FloodNet train$+$val, 2{,}057 images of which 855 are water-free; Section~\ref{sec:results-combined}), mean $\pm$ std over three folds. FloodNet test (424 images, 215 water-free) and AIFloodSense are never used for training or model selection in either regime; AIFloodSense has 94 images on the left and 93 on the right, after the corpus's duplicate filter removes one copy. Per-image macro metrics as in Table~\ref{tab:t9-zeroshot}; FAR is the fraction of water-free scenes with any predicted water (lower is better). All models FP32.}
\label{tab:t10-combined}
\small
\setlength{\tabcolsep}{2pt}
\begin{tabular}{llcccc@{\hskip 6pt}cccc}
\toprule
& & \multicolumn{4}{c}{\textbf{FSSD-trained (zero-shot)}} & \multicolumn{4}{c}{\textbf{Combined-trained}} \\
\cmidrule(lr){3-6}\cmidrule(lr){7-10}
& & \textbf{AIFS} & \multicolumn{3}{c}{\textbf{FloodNet test}} & \textbf{AIFS} & \multicolumn{3}{c}{\textbf{FloodNet test}} \\
\cmidrule(lr){4-6}\cmidrule(lr){8-10}
\textbf{Model} & \textbf{KD} & \textbf{IoU} & \textbf{IoU (all)} & \textbf{IoU (flood)} & \textbf{FAR} & \textbf{IoU} & \textbf{IoU (all)} & \textbf{IoU (flood)} & \textbf{FAR} \\
\midrule
Teacher (UNet$+$EffB0) & --   & 0.706 & 0.288 & 0.509 & 0.927 & 0.731 $\pm$ 0.011 & 0.783 $\pm$ 0.082 & 0.781 $\pm$ 0.003 & 0.216 $\pm$ 0.165 \\
\midrule
MobileNetV3-Small      & none & 0.698 & 0.285 & 0.506 & 0.929 & 0.713 $\pm$ 0.008 & 0.860 $\pm$ 0.023 & 0.754 $\pm$ 0.011 & 0.037 $\pm$ 0.036 \\
                       & comb & 0.693 & 0.280 & 0.497 & 0.930 & 0.725 $\pm$ 0.012 & 0.848 $\pm$ 0.018 & 0.780 $\pm$ 0.014 & 0.085 $\pm$ 0.047 \\
EfficientNet-Lite0     & none & 0.702 & 0.388 & 0.507 & 0.727 & 0.731 $\pm$ 0.005 & 0.854 $\pm$ 0.006 & 0.736 $\pm$ 0.015 & 0.031 $\pm$ 0.010 \\
                       & comb & 0.690 & 0.399 & 0.503 & 0.702 & 0.739 $\pm$ 0.008 & 0.844 $\pm$ 0.031 & 0.779 $\pm$ 0.015 & 0.093 $\pm$ 0.060 \\
MobileViT-XXS          & none & 0.723 & 0.253 & 0.512 & 1.000 & 0.738 $\pm$ 0.012 & 0.831 $\pm$ 0.017 & 0.695 $\pm$ 0.027 & 0.037 $\pm$ 0.009 \\
                       & comb & 0.708 & 0.262 & 0.507 & 0.975 & 0.727 $\pm$ 0.006 & 0.854 $\pm$ 0.011 & 0.762 $\pm$ 0.018 & 0.057 $\pm$ 0.035 \\
\bottomrule
\end{tabular}
\end{table*}


\emph{First, the false-alarm failure disappears.} Combined-trained students false-alarm on 3--9\% of the 215 water-free FloodNet scenes (Table~\ref{tab:t10-combined}), against 70--100\% for the same configurations trained on FSSD. Flood-only IoU on FloodNet rises from 0.50--0.51 to 0.70--0.78, and macro IoU from 0.25--0.40 to 0.83--0.86. This confirms that the failure in Section~\ref{sec:results-zeroshot} was a training-data property, not an architectural one: none of the three lightweight backbones lacks the capacity to reject a dry scene once dry scenes are supervised. AIFloodSense, which neither regime trains on, improves by a smaller margin (0.69--0.72 $\rightarrow$ 0.71--0.74).

\emph{Second, the teacher is the least reliable rejector.} Teacher false-alarm rates vary by more than an order of magnitude across folds (0.279, 0.028, 0.340), while every no-KD student stays at or below 0.08 in every fold. The teacher still has the highest flood-only IoU (0.781), so its lower macro IoU (0.783 vs.\ 0.831--0.860 for the students) comes entirely from false alarms. A plausible cause is model selection: the validation fold contains no water-free image, so checkpoint selection rewards flood segmentation and is blind to rejection behaviour. We have not tested this explanation.

\emph{Third, distillation now helps segmentation but transmits the teacher's false alarms.} On this benchmark combined KD raises FloodNet flood-only IoU for every student in every fold (9 of 9 student--fold pairs). Pooling per-image differences over the 209 flood-containing test images (averaged over folds, 10{,}000 bootstrap replicates), the gain is $+0.026$ [95\% CI $+0.015$, $+0.038$] for MobileNetV3-Small, $+0.043$ [$+0.030$, $+0.057$] for EfficientNet-Lite0, and $+0.067$ [$+0.048$, $+0.087$] for MobileViT-XXS. The same distillation raises the false-alarm rate by $+0.048$ [$+0.028$, $+0.070$] and $+0.062$ [$+0.040$, $+0.084$] for the first two students; for MobileViT-XXS the increase ($+0.020$ [$-0.000$, $+0.042$]) is not significant. The extra false alarms are concentrated on the scenes where the teacher itself false-alarms. Over all water-free test images and folds, the teacher false-alarms in 139 of 645 image--fold cases. On those, distilled students false-alarm 23--33\% of the time against 10--14\% for their no-KD counterparts; on the remaining 506 the rates are 0.6--3.2\% with or without distillation. The no-KD rate is itself elevated on the teacher's false-alarm scenes, so some of these images are intrinsically hard. To test whether the KD effect is specific to them, we compute a difference-in-differences: the KD-induced change in false-alarm rate on teacher-false-alarm cases minus the change on teacher-correct cases, with a bootstrap that resamples images (keeping each image's three fold copies together; 10{,}000 replicates; \texttt{scripts/stats\_combined.py}). It is $+0.186$ [$+0.114$, $+0.263$] for MobileNetV3-Small, $+0.169$ [$+0.093$, $+0.246$] for EfficientNet-Lite0, and $+0.139$ [$+0.069$, $+0.214$] for MobileViT-XXS. Distillation therefore adds false alarms mainly where the teacher makes them, although for EfficientNet-Lite0 it also adds a small number elsewhere ($+0.026$ on teacher-correct cases). This pattern is consistent with the student imitating the teacher's outputs on those scenes; our design does not separate the response and feature terms, so we cannot say which carries the effect. On AIFloodSense the effect of distillation is mixed ($+0.012$, $+0.008$, $-0.011$ IoU).

The negative-KD finding of Section~\ref{sec:results-kd} is therefore specific to the saturated FSSD regime. When the teacher has headroom over the students on the positive class (teacher flood-only IoU 0.781 vs.\ 0.695--0.754 without KD), distillation recovers most of that gap, but it imports the teacher's rejection errors along with it. For unattended monitoring, where false alarms carry the operational cost (Section~\ref{sec:discussion-asymmetry}), the no-KD students are the safer choice. For event-time mapping, where missed flood pixels matter more, combined KD is worth its cost only if the teacher's own false-alarm rate is first verified on water-free data.
